\section{Precisione}

\subsection{Macchine di Turing precise}
Una macchina di Turing $M$ é \textbf{precisa} se esistono due funzioni $f$ e $g$ tali che:
\begin{itemize}
	\item Ogni computazione di $M$ termina \textit{esattamente} dopo $f(n)$ passi.
	\item Al termine di ogni computazione tutti i nastri di $M$ sono lunghi \textit{esattamente} $g(n)$ (eccetto primo e ultimo nastro se la macchina é con input e output).
\end{itemize}

\subsection{Funzioni di complessitá proprie}
Le funzioni di complessitá proprie sono funzioni sui naturali, non decrescenti e tali che una macchina di Turing con input e output $M_f$ puó calcolarne il valore in unario in base alla lunghezza dell'input. \\
Cioé, dato un input $x$ lungo $n$, $M_f$ scrive sull'ultimo nastro tanti simboli $\sqcap$ ("quasi-blank") quanto vale $f(n)$. \\
Il tempo impiegato deve essere $O(n + f(n))$\footnote{Vedremo che $O(n + f(n)) = O(f(n))$}. \\
Lo spazio su ogni nastro deve essere $O(f(n))$. \\
Formalmente: una funzione $f: \mathbb{N} \rightarrow \mathbb{N}$ é \textit{propria} sse:
\begin{itemize}
	\item $f$ non é decrescente: \quad $\forall n \in \mathbb{N} \quad f(n+1) > f(n)$
	\item Esiste una MdT con input e output $M_f$ tale che, per ogni $n$: \\
	\[
	\configurazionestart{x} \xrightarrow{{M_f}^t} \configurazionepropria{x}{f(n)}
	\]
	\item $t = O(n + f(n))$
	\item $z_i = O(f(n)) \quad \forall i \in \{2, \dots, k-1\}$
\end{itemize}

Le funzioni proprie sono le uniche funzioni che possono essere usate per misurare la complessitá di tempo e di spazio di un problema senza incorrere in problemi di misurazione.\\ 
Ad esempio, il Gap Theorem dimostra che esistono funzioni che "rompono" l'analisi di complessitá.

\subsubsection*{Gap Theorem}
Esiste una funzione ricorsiva $f: \mathbb{N} \rightarrow \mathbb{N}$ tale che la classe di problemi risolvibili in tempo $f(n)$ coincide con quella di problemi risolvibili in tempo $2^{f(n)}$. \footnote{Vedremo meglio l'enunciato una volta introdotta la classe di complessitá \TIME{}}
